\section{The actual level cover and its formal theta connection}
\label{sec:actual-level-cover}

The local integral charts already constructed do not by themselves
provide a connection over every coefficient algebra.  The latter is a
separate quantifier in the original level-cover datum.  Over the original
coefficient field it can be supplied directly by the formal exponential.
This also constructs the actual finite cover with its framed-point
lifting property.  No new number field is required.

Write $K=c.F$, $\mathcal O=\mathcal O_E$, and retain the original
datum, primitive cubical line, and coefficient embeddings.  In
particular $K$ contains all conjugates of $E$ and the prescribed
cyclotomic roots, and the original rank satisfies $p+q\ge1$.
All moduli schemes below are the actual ones already
constructed in Sections~\ref{sec:integral-pel-frontier} and
\ref{sec:pel-level-system}.

\begin{theorem}\label{thm:actual-level-cover}
The following suppliers exist with the original source contracts.
\begin{enumerate}
\item On each original moduli scheme there is a family of formal theta
derivatives satisfying \texttt{IsNab}, for every commutative $K$-algebra
$S$, including nonreduced and nonnoetherian algebras.  Every original
framed point over a reduced $K$-algebra satisfies \texttt{FromTest}.
The original harmonic combinations of these derivatives commute with
framed base change over arbitrary $K$-algebras.
\item Suppose $l\mid l_2$, $M_2=M p_d$, $(p_d,M)=1$,
$\mathfrak b_2=\mathfrak b I$, and $I+\bar I=\mathcal O$.
Suppose the original coefficient data have the same field
$c_2.F=c.F$.  For every $n\ge0$ there is an actual
$d:\texttt{DeltaDat}(X,c,n)$ with precisely these $p_d,I,l_2,X_2,c_2$,
with $c_\delta=1$, $\phi_\delta=1$, and
\[
               d.\texttt{IsLev},\qquad
               \texttt{IsNab}(X_2,c_2,d.\nabla_2).
\]
The morphism on generic moduli over the common field is finite,
\'etale and surjective.  Its map of universal families is the original
level-forgetting isomorphism and preserves the primitive line,
rigidification and cube structure.  Every original framed geometric
point lifts with its exact tangent frames and section generator.
\end{enumerate}
\end{theorem}

\begin{proof}
\paragraph{Formal coordinates over an arbitrary coefficient algebra.}
Let $y$ be an original framed point over $S$ and let
$A/S$ be its abelian scheme.  The finite \'etale algebra
\[
                   E\otimes_{\mathbb Q}K
                       \simeq\prod_{\tau:E\hookrightarrow K}K
\]
splits after tensoring with $S$.  Its orthogonal idempotents decompose
$\operatorname{Lie}(A/S)$ into exactly the eigenspaces in the source
definition of $y$.  Thus the specified separate eigenbases concatenate
to a basis $v_1,\ldots,v_g$ of the full tangent module.  This argument
uses units in $K$; it makes no assertion about the sum of the separate
integral eigenspaces over a valuation ring.

Translate $v_i$ to the invariant derivations $D_i$ on the formal group.
They commute because $A$ is commutative.  For a formal function $f$ set
\begin{equation}\label{eq:level-cover-formal-exponential}
       \exp_A^*f
          =\sum_{\alpha\in\mathbb N^g}
             \frac{(D^\alpha f)(0)}{\alpha!}\,t^\alpha.
\end{equation}
Each coefficient uses only a finite derivative; the series is formal.
The Leibniz formula proves that this is a homomorphism of complete
$S$-algebras.  Invariance and commutation prove that it carries group
addition to $t+t'$.  Its linear term is the identity on the specified
cotangent basis.  The ideal of a smooth section is a regular ideal
with associated graded algebra the symmetric algebra of its conormal
module.  Induction on the infinitesimal order consequently makes
\eqref{eq:level-cover-formal-exponential} an isomorphism at every finite
order and then on the inverse limit.  This proof is valid over any
$\mathbb Q$-algebra; neither noetherianity nor reducedness is used.
It also shows directly that the original coefficient action becomes
its given linear tangent action.  The construction commutes with
arbitrary change of $K$-algebra.

\paragraph{The original biextension and primitive line.}
Pull the actual rigidified biextension to
$\operatorname{Spf}S[[t,z]]$.  It has a frame compatible with its
zero rigidification.  Indeed, at each finite order the coordinate ring
is a nilpotent extension of $S$.  An invertible module whose reduction
has a chosen generator has a lift of that generator, and the lift is a
generator by the nilpotent form of Nakayama's lemma.  Successive
generators can be made compatible: their ratio on the preceding order
is a unit, which lifts across a nilpotent ideal.  Passing to the inverse
limit gives a formal frame.  If its two axis values are $a(t)$ and
$b(z)$, multiplication by $a(0)/(a(t)b(z))$ gives the exact prescribed
frames on both axes.

Use the transported \emph{original} primitive identity on the graph
$(\mathrm{id},w_0)$ to frame the original line $L$; this is not a new
theta characteristic or a choice made merely from its polarization.
The actual cube, symmetry and adjoint identities give the four formal
cocycles used in \texttt{ChartAt}.  Their constant terms are one.
The gauge construction in the proof of Theorem~\ref{thm:actual-good3}
works verbatim over $S$: its Euler homotopies divide only by positive
integers, its adjoint correction uses the separability idempotent of
$E/\mathbb Q$, and its symmetry correction divides by two.  All these
scalars are units in $S$.  No step uses division by a nonzero element
of $S$ or passage to a fraction field.  Since the coordinates above
already linearize the group and coefficient action, this constructs
the required canonical gauge for the actual framed biextension.

Every actual section of $L^k$, $k\ge0$, has its formal germ in this
frame: restrict it to each infinitesimal neighbourhood and use the
compatible trivializations.  The restrictions are compatible before
taking the limit.  Therefore all the original \texttt{FrTruncC},
\texttt{IsPrim}, and \texttt{IsGerm} clauses hold, not just equality
of values at the zero section.  Choose such a chart for each $S,y$
and define its theta derivatives by the source's ordinary-derivative
formula
\[
     \nabla_y^\alpha(F)
       =\left.\partial^\alpha
          \bigl(\exp(kh(t,w_0t))\,f_F(t)\bigr)\right|_{t=0},
                    \qquad k\ge0.
\]
If the source regards the power as an arbitrary integer, set the
unused negative-power entries to zero, as its \texttt{IsNab} contract
does.  Uniqueness of the formal germ proves that this definition
satisfies the entire \texttt{IsNab} clause, including its universal
quantifier over all representatives of the germ.

The choice of gauge need not give canonical nonharmonic derivatives.
For harmonic combinations it does: two gauges differ by a balanced
symmetric bilinear form $B$ in the additive coordinates, and
$B(t,w_0t)=B(t,t)/2$.  The associated finite differential correction
is a sum of the original paired Laplacians and vanishes on the
original harmonic polynomial.  This proof is a coefficient identity
over $S$.  Distinct coefficient embeddings remain separated by units
coming from $K$, so it is still valid with nilpotents and after
nonflat base change.  Comparing a chosen target chart with the
base-changed source chart proves the asserted harmonic compatibility.

\paragraph{Framed points and the actual test rings.}
On the moduli scheme over $K$, the eigenspaces of the Hodge bundle
are vector bundles of the prescribed ranks.  The actual theta
pushforward is an invertible sheaf and commutes with base change,
by the previously constructed section supplier, restricted to the
abelian locus.  Thus the framed-point functor is represented by the
product of their frame bundles and the nonzero-vector bundle of
this theta line.  It is a scheme of finite presentation over $K$.

Write an arbitrary reduced $K$-algebra $S$ as the filtered union of
its finitely generated $K$-subalgebras.  A morphism from
$\operatorname{Spec}S$ to this finite-presentation frame scheme
descends to some such subalgebra $S_0$; this follows from the
finite-presentation limit theorem \cite{level-cover-limits}.
The algebra $S_0$ is reduced because it is a subalgebra of $S$.
Its point is an actual framed point and pulls back to $y$, with
its section and every tangent vector.  Presenting $S_0$ as a quotient
of a polynomial algebra gives precisely an original reduced test
ring.  This proves \texttt{AllTest}, including every algebraically
closed coefficient field.  In particular the natural rule $1$ has
value exactly one at all the framed points used below.

\paragraph{Forgetting the added level.}
Theorem~\ref{thm:original-half-level-completion} identifies the
original moduli problem over $K$ with its full, symplectically
liftable $\operatorname{lcm}(M,A)$-marking, using the original
cyclotomic root.  Apply it at both levels.  The completion preserves
the original labels even when $M$ and $A$ share prime powers.
Forgetting the extra finite Tate labels is represented by a finite
\'etale cover: the possible compatible markings form the finite
locally constant sheaf of reductions of integral Tate isometries.
At each geometric point an original full Tate marking supplies a
lift, so the cover is surjective.  Here liftability is part of the
original marking condition; an arbitrary finite symplectic marking
is not substituted.  The simultaneous root adjustment in the
half-level completion theorem preserves all the old labels and
chooses the prescribed new roots.  More explicitly, put
$L=\operatorname{lcm}(M,A)$ and $L_2=\operatorname{lcm}(M_2,A_2)$.
Starting with a full Tate marking having the prescribed $L$-root,
the multiplier needed for the prescribed $L_2$-root is
$r_0\in\widehat{\mathbb Z}^{\times}$ with $r_0\equiv1\pmod L$.
At every finite place the original extension $E/E^+$ is split or
unramified.  Norm surjectivity on the corresponding principal-unit
groups gives $u\in\widehat{\mathcal O}^{\times}$ with
$u\bar u=r_0$ and $u\equiv1\pmod{L\widehat{\mathcal O}}$.
Composing with $u$ changes the cyclotomic multiplier and preserves
the entire old $L$-marking.  Thus the prescribed higher-root
condition does not remove any geometric point of the cover.

Before extension to $K$, the same forgetful natural transformation
gives a morphism $r:X_2\to X$ over the prescribed map of cyclotomic
base fields.  Its universal abelian scheme is the pullback of the
old one, with its coefficient action and polarization.  Functoriality
of the Poincar\'e bundle and the original primitive graph gives the
exact isomorphism of rigidified cubical lines and hence the source
map $m$ of universal semiabelian data on these generic moduli.
The two label identities are exactly
\begin{equation}\label{eq:level-cover-labels}
 m\circ[p_d]\eta_{M_2}(a)=\eta_M(a)\circ r,
 \qquad
 m\circ[A_2/A]\eta_{A_2}(a)=\eta_A(a)\circ r.
\end{equation}
They follow by restricting the same Tate marking: a point labelled
$a/M_2$ becomes $a/M$ after multiplication by $p_d$, and similarly
for $A$.  The forced complementary labels agree by the perfect
pairing, so these equations hold for the entire original maps,
not only their selected half-submodules.

Let $y$ be a framed point over any algebraically closed $K$-field.
Choose a geometric point $x_2$ above it, using surjectivity.  Pull
back its tangent bases and its exact section generator through the
universal-family and primitive-line isomorphisms.  These are actual
isomorphisms, so the tangent eigenspace bases and rank-one section
module remain exactly the ones specified by $y$.  This gives the
full \texttt{FrOver} identity, including the local transition-frame
formula for the section.  The construction also commutes with base
change and triangular changes of the prescribed tangent frames.

Finally choose the connection constructed above on $X_2$ and set
$c_\delta=1$ and $\phi_\delta=1$.  The constant rule is a section
of weight $(0,0)$, and its value on every framed geometric point
is one by \texttt{AllTest}.  The arithmetic level hypotheses,
\eqref{eq:level-cover-labels}, framed lifting, and this nonvanishing
are exactly the successive fields of \texttt{DeltaDat.IsLev}
\cite{level-cover-source}.  This proves both assertions.
\end{proof}

\paragraph{What this removes.}
The level cover, connection, test-ring descent and framed-point lift
are now actual suppliers on the original arithmetic moduli.  This
result introduces no component-field extension or additional local
lattice index.  It does not yet construct the full sized
\texttt{Deep.Good} cover.  Its remaining inputs include the exact
component-field statement and the complex \texttt{Member}
uniformization with its probability-normalized theta section and
squared multiplier identity.  The theta-series/section identification
and its global norm estimates remain separate obligations.
